\section{Fundamental Theorem of Calculus}

\begin{theorem}[Mean Value Theorem]%
  \label{th:media_integrale}%
  \mymark{***}%
  Let $a,b\in \RR$, $a<b$, and let
  $f\colon \closeinterval{a}{b} \to \RR$ be a Riemann-integrable function.
  Then
  \begin{equation}\label{eq:media_inf_sup}
    \inf_{\closeinterval{a}{b}} f 
    \le \frac{\int_a^b f(x)\, dx}{b-a} 
    \le \sup_{\closeinterval{a}{b}} f.
  \end{equation}
  Furthermore, if $f$ is continuous, 
  there exists a point $c \in \closeinterval{a}{b}$
  such that
  \[
  \frac{\int_a^b f}{b-a} = f(c).
  \]
  \end{theorem}
  %
  The quantity
  \[
    \frac{\int_a^b f}{b-a}
  \]
  is called the \emph{integral mean value} of $f$ on $[a,b]$ and is often
  denoted by the symbol
  \[
    \dashint_a^b f.
  \]
  %
  \begin{proof}
  \mymark{***}
  For every $x\in [a,b]$, we have 
  \[
    \inf_{\closeinterval{a}{b}} f \le f(x) \le \sup_{\closeinterval{a}{b}}f.
  \]
  Integrating, we obtain
  \[
    \inf_{\closeinterval{a}{b}} f \cdot (b-a)
    \le \int_a^b f(x)\, dx 
    \le \sup_{\closeinterval{a}{b}}f\cdot (b-a)
  \]
  and dividing by $b-a$, we arrive at~\eqref{eq:media_inf_sup}.

  If the function $f$ is continuous on $[a,b]$, by the Weierstrass theorem, 
  $\inf$ and $\sup$ are values assumed by the function, 
  and the integral mean $\dashint_a^b f$ is an intermediate value between the 
  minimum and the maximum.
  But by the intermediate value theorem~\ref{th:valori_intermedi}, 
  we can finally assert that there exists a point $c\in [a,b]$
  where the function assumes this value.
\end{proof}
  
\mymargin{fundamental theorem}
\begin{theorem}[Torricelli-Barrow: Fundamental Theorem of Calculus]
\mymark{***}%
\index{theorem!fundamental theorem of calculus}%
\index{theorem!Torricelli-Barrow}%
\index{Torricelli-Barrow!theorem of}%
\label{th:torricelli-barrow}%
Let $I\subset \RR$ be an interval, let $x_0 \in I$,
 and let $f\colon I\to \RR$ be a continuous function.
Then the \emph{integral function}
\index{function!integral function}
$F\colon I \to \RR$
\mymargin{integral function}
\[
  F(x) = \int_{x_0}^x f
\]
is well-defined, is differentiable, and for every $x\in I$, we have
\[
  F'(x) = f(x).
\]
In particular, since $f\in C^0(I)$, we have $F\in C^1(I)$.

Moreover, if $G\colon I \to \RR$ is any function such that
$G'(x) = f(x)$ for every $x\in I$, then for every $a,b \in I$, we have
\mymargin{fundamental formula of calculus}
\index{formula!fundamental formula of calculus}%
\[
  \int_a^b f = G(b) - G(a).
\]
\end{theorem}
%
\begin{proof}
\mymark{***}
First, observe that the function $f$, being continuous, is integrable
on every closed and bounded interval contained in $I$. Thus, the integral
$\int_{x_0}^x f$ is well-defined.

For every $h\neq 0$, if $x+h \in I$, by the additivity of the integral,
we have
\[
\frac{F(x+h) - F(x)}{h} = \frac{\int_{x_0}^{x+h} f - \int_{x_0}^x f}{h}
 = \frac{\int_x^{x+h} f}{h}.
\]
Applying now the mean value theorem, we can
assert that there exists a point $\xi(h)$ in the interval with endpoints $x$ and
$x+h$
such that
\[
  \frac{\int_x^{x+h} f}{h} = f(\xi(h)).
\]
As $h\to 0$, we have $\xi(h) \to x$ and, by the continuity of $f$,
$f(\xi(h)) \to f(x)$.
Thus, we have shown that $F$ is differentiable at $x$:
\[
 \lim_{h\to 0}\frac{F(x+h)-F(x)}{h} = f(x)
\]
and $F'(x) = f(x)$.

Thus, if $a,b\in I$ are any points, we have:
\[
\int_a^b f = \int_{x_0}^b f - \int_{x_0}^a f = F(b) - F(a).
\]
And if $G\colon I \to \RR$ is any function such that $G'(x)=f(x)$, we
will have $G'(x) = F'(x)$ for every $x\in I$, and thus $(G-F)' = 0$ on $I$.
By the monotonicity criteria, we can conclude that $G-F$ is constant on $I$:
$G-F = c$. Thus, we have
\[
 \int_a^b f = F(b) - F(a) = (G(b) - c) - (G(a) - c) = G(b) - G(a).
\]
\end{proof}

\begin{example}
We wish to calculate
\[
  \int_0^b x^2\, dx.
\]
It will suffice to observe that setting $G(x)=\frac{x^3}{3}$, we have
$G'(x)=x^2$,
and thus, thanks to the fundamental formula of calculus, we have
\[
  \int_0^b x^2\, dx = G(b) - G(0) = \frac{b^3}{3}.
\]
It is evident how this method of resolution is much simpler
and more powerful than the one used in example~\ref{ex:integrale_quadrato}.
\end{example}

\begin{definition}[Primitive]
\label{def:primitiva}%
\mymark{***}%
Let $A \subset \RR$ and let $f\colon A \to \RR$ be any function.
A function $F\colon A \to \RR$ is said to be a \emph{primitive}%
\mymargin{primitive}%
\index{primitive}
(or \emph{antiderivative})
\index{antiderivative}
of $f$ if $F$ is differentiable and $F'(x)=f(x)$ for every $x\in A$.
\end{definition}

The fundamental theorem of calculus can thus be expressed in
the following way: every continuous function $f$, defined on an interval,
admits at least one primitive, and if $F$ is any primitive of $f$, we have
\[
  \int_a^b f = F(b) - F(a).
\]
To denote the difference $F(b)-F(a)$, the following notations are sometimes
used:
\[
  \Enclose{F(x)}_{x=a}^b = \Enclose F_a^b
  = F(x) \vert_{x=a}^b
  = F \vert_a^b = F(b) - F(a).
\]

The calculation of integrals is thus reduced to the determination of primitives,
i.e., to inverting the derivative operator.
It will therefore be important to have tools for determining the primitives
of a function.

\begin{definition}[Indefinite Integral]
\mymark{***}
The set of all primitives of a function $f\colon A \to \RR$
is denoted by the symbol
\[
  \int f
  \qquad\text{or}\qquad
  \int f(x) \, dx
\]
and is called the \emph{indefinite integral}.
The reason for this notation (and the name) derives from the fundamental theorem
of
calculus, according to which if $f\colon[a,b]\to \RR$
is continuous, we have
\[
  \int_a^b f = \Enclose{\int f}_a^b.
\]
\end{definition}

\begin{remark}
Note, however, that in the case of non-continuous functions, it is possible
that the integral functions are not primitives. For example,
if we take the Heaviside function $H(x)$ defined in example~\ref{ex:heaviside},
we can observe that $F(x)=\int_0^x H(t)\, dt = \abs{x}$
is $0$ if $x \le 0$ and is $x$ if $x \ge 0$.
Thus, $F$ is an integral function but not a primitive of $H$
(and the set of primitives, in this case, is empty).
\end{remark}

If we think of the linear operator $D$ defined on the set of differentiable
functions
$Df = f'$,
we can think of $\int f$ as the set of preimages of $f$
via $D$, i.e.:
\[
  \int f = D^{-1}(\ENCLOSE{f}) = \ENCLOSE{F \colon DF =f }.
\]
Thus, $\int f$ is the set of solutions $F$ of the non-homogeneous linear
equation
$DF =f$. The associated homogeneous equation is $DF = 0$.
If the functions are defined on an interval, then the solutions of the
homogeneous equation are constants (as stated in the following theorem).
Thus, in such a case, the space $\int f$ is a one-dimensional affine space
parallel to the vector space of constant functions. 
But if we choose a domain that is not an interval 
(as in example~\ref{ex:primitive_non_connesso}), 
we will have a space whose dimension is equal to the number of intervals that
compose the domain.

\begin{theorem}[Properties of Primitives]
\label{th:primitive}
\mymark{***}
Let $f\colon I \to \RR$
be a continuous function defined on a non-empty interval $I\subset \RR$. Then
\begin{enumerate}
\item there exists at least one primitive $F$ of $f$;
\item given a primitive $F$ of $f$, every other
primitive $G$ differs from $F$ by a constant: $\exists c\in \RR\colon G= F+c$.
\end{enumerate}

In other words, if $f$ is continuous, $\int f$ is not empty, and if 
the domain of $f$ is an interval and $F\in \int f$ is any primitive, then
\[
  \int f = \ENCLOSE{F+c \colon c \in \RR}.
\]
\end{theorem}

Observe that the set of constant functions
on an interval
is nothing but $\ker D$, i.e., the null space of the derivative operator.
We are thus simply observing that the preimages of a linear operator
are affine spaces parallel to the kernel of the operator.

\begin{proof}
\mymark{***}
Choosing a point $x_0\in I$, we can
consider the integral function
\[
  F(x) = \int_{x_0}^x f(t)\, dt.
\]
The fundamental theorem of calculus
assures us that $F$ is a primitive of $f$.

Conversely, if $F$ and $G$ are two primitives of $f$, then we have:
\[
  F' = G' = f.
\]
Setting $H=G-F$, we will have $H'=0$ on the interval $I$. By the monotonicity
criteria, we know that $H$ is constant, i.e., there exists $c\in \RR$ such
that $H(x)=c$ for every $x\in I$. Thus, we obtain, as desired: $G=F+H=F+c$.
\end{proof}

If the function $f$ is not defined on an interval, the set of primitives
will have a dimension greater than $1$, as seen in the following.

\begin{example}[Primitives on Non-Connected Sets]
\label{ex:primitive_non_connesso}%
\mymark{*}%
Consider the function $f(x) = 1/x$. Observe that
$f\colon (-\infty, 0) \cup (0,+\infty)\to \RR$ is defined on the union of two
intervals. By direct verification, we can observe that the function
$F(x) = \ln \abs{x}$ is a primitive of $f$. To obtain the set of all
primitives, we can add to $F$ any function with a zero derivative on the domain
of $f$. The functions with zero derivative are constant on each
interval, and thus we find that for every $c_1, c_2\in \RR$, the function
\[
G(x) =
\begin{cases}
  \ln (x) + c_1 &\text{if $x>0$},\\
  \ln (-x) + c_2 & \text{if $x<0$}
\end{cases}
\]
is a primitive of $f$, and there are no other primitives.

In this case, the space of primitives has dimension $2$ because the kernel
of the derivative operator on the space of functions defined on the union of
two intervals has dimension $2$.
This is the simplest example of a rather general phenomenon where differential
operators on a certain space are closely related to the topology of the space
itself. In this case, the dimension of the kernel
of the derivative operator $D$ is equal to the number of connected components of
the domain of the functions in the domain of $D$.
\end{example}